
\section{One Claim, Three Questions}

The claim under test is familiar: Peter entered the empty tomb, saw the linen
cloths arranged in a remarkable way, and immediately understood that Jesus had
risen. It often arrives with a second claim: the face-cloth was folded according
to an ancient Jewish table custom meaning that a master would return. The two
claims are rhetorically effective because they turn otherwise quiet objects
into a coded message.

The canonical text is both firmer and more restrained. John explicitly says
that Peter entered and contemplated the cloths. John explicitly says that the
other disciple entered, saw, and believed. He does not attach the verb
``believed'' to Peter. Luke says that Peter saw the cloths and departed
wondering. Neither account describes a table custom or calls the face-cloth a
napkin in the modern dining sense.

Three questions must therefore stay apart:

\begin{enumerate}
\item What did Peter observe, and what state of mind does each Evangelist
attribute to him?
\item What can the disposition of the cloths reasonably establish about the
absence of the body?
\item Is the folded-napkin return signal an ancient interpretation, or a modern
story attached later to the passage?
\end{enumerate}

The resulting judgment is divided. The cloths are real narrative evidence, and
ancient Christian interpreters used their disposition against theft. John,
however, attributes belief to the other disciple, while Luke leaves Peter in
wonder. The modern table-signal explanation is unsupported by the named Jewish,
patristic, and historical witnesses searched for this study.

